\note{How do developers test protocols today, and why isn't that enough?\par \par First, \textbf{hand-designed tests} run a protocol under manually chosen parameters, such as fixed bandwidth and latency ranges. They are easy to interpret and great for checking a suspected issue, but their coverage is limited to scenarios the designer already thought of testing—they cannot cover the exponential space of time-varying sequences.\par \par Second, \textbf{real-world trace replay and testbeds}—such as replaying captured traces in Mahimahi, or evaluating protocols on platforms like Pantheon and Puffer—test conditions that users actually experience. These are essential for benchmarking average real-world performance, but their coverage is bounded by where and when traces were collected. They do not actively search for unknown worst-case conditions, and when an anomaly does happen in live deployment, uncontrolled background processes make it very hard to reproduce.
\par\smallskip Source: AdvNet paper §1–3}
